As it can be noted in its description the SCE is easly applied to any
optimization problem.
The only steps needed to be specified is the creation of a new random
solution (Algorithm 1) and the crossing procedure of two solutions (Algorithm 2).
%These two procedures for the MKP are respectively presented by the algorithms
%below.
\begin{center}
\includegraphics[width=0.95\linewidth]{imgs/alg-new}\\[3mm]
%To construct a new random solution (Algorithm~\ref{alg:new}) the items are
%at first shuffled in random order and stored in a list (line 2).
%A new empty solution is then defined (line 3).
%The algorithm iteratively tries to fill the solution's knapsack with 
%the an item taken from the list (lines 4-9).
%The feasibility of the solution is then checked: if the item insertion let
%the solution unfeasible (line 6) its removed from knapsack (line 7).
%After trying to place all available items the new solution is returned.
\includegraphics[width=0.95\linewidth]{imgs/alg-cross}
\end{center}
%\vspace{0.1mm}
